% Chapitre 2 - Page 9 | Travaux Avant-Projet | L3 ELN
\documentclass[11pt,a4paper]{article}
\usepackage{fontspec}
\setmainfont{TeXGyrePagella}
\usepackage[french]{babel}
\usepackage{graphicx,tikz,xcolor,geometry}
\usetikzlibrary{calc}
\geometry{margin=0pt}
\pagestyle{empty}
\definecolor{bleufonce}{HTML}{1B3A6B}
\definecolor{bleugris}{HTML}{5C7190}
\newcommand{\pos}[4]{\node[anchor=north west,inner sep=0pt,outer sep=0pt,text width=#3,align=left] at ($(current page.north west)+(#1,-#2)$) {#4};}
\newcommand{\sect}[1]{{\color{bleufonce}\bfseries\fontsize{17}{19}\selectfont #1}}
\newcommand{\step}[1]{{\color{bleufonce}\bfseries\fontsize{12.5}{14.5}\selectfont #1}}
\newcommand{\smallfont}{\fontsize{9.8}{12}\selectfont}
\newcommand{\legende}[1]{{\color{bleugris}\fontsize{8.5}{10}\selectfont #1}}
% \shot{x}{y}{largeur}{hauteur}{fichier} : capture encadrée
\newcommand{\shot}[5]{\node[anchor=north west,inner sep=0pt,outer sep=0pt] at ($(current page.north west)+(#1,-#2)$) {\includegraphics[width=#3]{#5}};
  \draw[bleufonce!75,line width=.55pt,rounded corners=1mm]
    ($(current page.north west)+(#1,-#2)+(-.4mm,.4mm)$) rectangle ($(current page.north west)+(#1,-#2)+(#3,-#4)+(.4mm,-.4mm)$);}
\begin{document}
\begin{tikzpicture}[remember picture,overlay]
\draw[bleufonce,line width=.95pt,rounded corners=8mm]
 ($(current page.north west)+(5mm,-5mm)$) rectangle ($(current page.south east)+(-5mm,5mm)$);

\pos{10mm}{13mm}{189mm}{\sect{6. Créer un montage dans Fritzing}}
\pos{10mm}{25mm}{188mm}{{\fontsize{10.5}{13}\selectfont Les étapes ci-dessous sont illustrées par des captures de \textbf{Fritzing 1.0}. L'interface étant en anglais, le nom de chaque commande est donné tel qu'il apparaît à l'écran.}}

% --- Étape 1 ---------------------------------------------------------------------
\pos{10mm}{40mm}{188mm}{\step{Étape 1 --- Ouvrir Fritzing}}
\shot{14mm}{49mm}{104mm}{55.1mm}{assets/F02_accueil.png}
\pos{125mm}{50mm}{72mm}{{\smallfont Au lancement, Fritzing s'ouvre sur l'onglet \textbf{Welcome} (accueil) :\\[1.8mm]
  \textcolor{bleufonce}{\textbullet}\enspace \textbf{Recent Sketches} : les derniers croquis ouverts ;\\[1.2mm]
  \textcolor{bleufonce}{\textbullet}\enspace \textbf{New Sketch} / \textbf{Open Sketch} : créer ou ouvrir un croquis ;\\[1.2mm]
  \textcolor{bleufonce}{\textbullet}\enspace à droite : actualités (\textbf{Blog}) et service de fabrication de PCB (\textbf{Fab}).}}
\pos{14mm}{106mm}{184mm}{\legende{Figure 7 --- L'écran d'accueil (\textbf{Welcome}) de Fritzing.}}

% --- Étapes 2 et 3 --------------------------------------------------------------
\pos{10mm}{116mm}{90mm}{\step{Étape 2 --- Nouveau croquis}}
\shot{14mm}{125mm}{54mm}{36.6mm}{assets/F03_nouveau_croquis.png}
\pos{14mm}{165mm}{84mm}{{\smallfont Cliquer sur \textbf{New Sketch} (nouveau croquis) ou utiliser \textbf{File \(\rightarrow\) New} (\textbf{Ctrl+N}). Un plan de travail avec une plaque d'essai s'ouvre.}}

\pos{108mm}{116mm}{90mm}{\step{Étape 3 --- Chercher un composant}}
\shot{112mm}{125mm}{44mm}{44.3mm}{assets/F04_recherche.png}
\pos{161mm}{126mm}{37mm}{{\fontsize{9.4}{11.5}\selectfont Dans le panneau \textbf{Parts}, cliquer sur la \textbf{loupe}, taper le nom en anglais (\emph{resistor}, \emph{LED}, \emph{photocell}\ldots) puis valider.}}
\pos{112mm}{172mm}{86mm}{{\smallfont Il suffit ensuite de \textbf{glisser} le composant vers le plan de travail.}}
\pos{14mm}{184mm}{184mm}{\legende{Figures 8 et 9 --- Création d'un croquis et recherche d'une résistance dans la bibliothèque.}}

% --- Étape 4 ---------------------------------------------------------------------
\pos{10mm}{192mm}{188mm}{\step{Étape 4 --- Placer et relier les composants}}
\shot{14mm}{201mm}{68mm}{54.7mm}{assets/F05_exemple_breadboard.png}
\pos{90mm}{202mm}{107mm}{{\smallfont
  \textcolor{bleufonce}{\textbullet}\enspace Placer l'\textbf{Arduino Uno} et les composants sur la plaque d'essai.\\[1.3mm]
  \textcolor{bleufonce}{\textbullet}\enspace Une patte bien insérée devient \textcolor{green!50!black}{\textbf{verte}} : elle est connectée.\\[1.3mm]
  \textcolor{bleufonce}{\textbullet}\enspace Pour tracer un fil : \textbf{cliquer-glisser} depuis une broche jusqu'à un trou.\\[1.3mm]
  \textcolor{bleufonce}{\textbullet}\enspace Un \textbf{clic droit} sur un fil permet de changer sa couleur (rouge : 5~V, noir : GND).\\[1.3mm]
  \textcolor{bleufonce}{\textbullet}\enspace La valeur d'une résistance se règle dans le panneau \textbf{Inspector} (menu \textbf{Window}).}}
\pos{14mm}{258mm}{184mm}{\legende{Figure 10 --- Notre exemple en version simplifiée : la photorésistance en pont diviseur sur \textbf{A0} et une \textbf{LED} commandée par la broche \textbf{2}.}}

\draw[bleufonce!55,line width=.45pt] ($(current page.south west)+(10mm,14.7mm)$)--($(current page.south east)+(-10mm,14.7mm)$);
\node[anchor=south west,inner sep=0pt,text=bleufonce!80] at ($(current page.south west)+(10mm,7.4mm)$) {\fontsize{9.5}{11}\selectfont 2026/2027};
\node[anchor=south east,inner sep=0pt,text=bleufonce!80] at ($(current page.south east)+(-10mm,7.4mm)$) {\fontsize{9.5}{11}\selectfont Page 9};
\end{tikzpicture}
\null
\end{document}
